\scanpage{6}
Pour le \unsure{composé|compos.} dans l'autre sens, on trouve que $G$ opère sur l'ens $I_G$ de ses trois ss-groupe de Sylow par autom. intérieurs, d'où
\[
G\xrightarrow{\sim}\unsure{\mathfrak{S}|G}_{I_G}
\]
qui est un iso.

Éléments de $\mathfrak{S}_I$.

a) $1$

b) Les \inserted{trois} éléments \deleted{non nuls de} \inserted{de} $V^{*}$, qui forment une orbite $J$ sous $\mathfrak{S}_I$, et correspondent aux élém. de $\mathcal{P}_{22}(I)$
% slanted
\margin{``translations'' (non nulles) [sym. p.r. aux trois axes du cube]}

c) Les \deleted{\unsure{deux}} $8$ éléments d'ordre $3$, se répartissent en \deleted{quatre} \inserted{deux} \deleted{torseurs} paquets de $4$ qui sont des torseurs (à gauche ou à droite, au choix) sous $V$, suivant les deux éléments $\ne1$ de $\mathfrak{S}_J^+$, L'ens de ces éléments s'identifie aussi à $I\times\underline{\omega}$, où $\underline{\omega}=\omega(\unsure{J|I})\bigl(\simeq\omega(I)\bigr)$, par l'application qui associe à tout $\unsure{\rho|p}\in\mathfrak{S}_I$ d'ordre $3$ son unique pt fixe dans $I$ (ou encore, le sous-groupe de Sylow qu'il engendre), \uline{et} l'orientation qu'il détermine \deleted{\unsure{via son}} ou encore son image dans $\mathfrak{S}_J^+$.
% slanted
\margin{``rotations'' (ayant un pt fixe et un seul, lequel transforme $I$ en vectoriel \unsure{isom.} à $V$) [``rotations-sommets'' du cube]}

d) Les \deleted{\unsure{12}} \inserted{$6$} éléments \deleted{\unsure{non}} \inserted{d'ordre 2} de $\mathfrak{S}_I$ qui s'envoient dans $\mathfrak{S}_J^-$, en correspondance 1-1 avec l'ens $\mathcal{P}_2(I)$ -- ce sont les \uline{transpositions} dans $I$ qui se groupent \deleted{par} \inserted{en 3} \deleted{\unsure{gr}} paires, pour tte transposition $\tau_A$, $\tau_{A'}$ \deleted{$\unsure{\tau}$} $(A'=I\setminus A)$ est l'unique \inserted{autre} élément d'ordre $2$ tel que $\tau_A\tau_{A'}\in V^*$.
% slanted
\margin{``réflexions'' (ayant chacune une \deleted{\unsure{hyperplan}} droite de pts fixes) [``rotations-arêtes'' du cube]}

e) Les $6$ éléments d'ordre $4$ de $\mathfrak{S}_I$, correspondent aux $6$ ordres circulaires sur $I$. Pour
% slanted
\margin{[``rotations-faces'' du cube]}
% archive sheet number: 5
